\section*{Problem 4: asymptotic PTAS for Bin Packing with trees}
\subsection*{The problem}
In the Bin Packing with Trees problem we are given a collection of $n
> 0$ items $I = \{1, 2, \ldots, n\}$, each item has a size $s_i$, $0
< s_i \leq 1$. Also, we are given an undirected tree $T = (V, E)$.
That is, the set of vertices of $T$ is the set of items (i.e., $V =
I$). A feasible solution for the problem is an m-tuple $U = (U_1,
U_2, \ldots, U_m)$ for some $m \geq 1$, such that
\begin{itemize}
  \item For every $j$, $1 \leq j \leq m$, $U_j \subseteq I$ and
    $\sum_{i \in U_j} s_i \leq 1$.
  \item For every $i \in I$ there is $j \in [m]$ such that $i \in U_j$.
  \item For all $j \in [m]$ and $\{u, v\} \in E$ it holds that $|U_j
    \cap \{u, v\}| \leq 1$. That is, two items connected by an edge
    cannot be packed in the same bin.
\end{itemize}
An asymptotic PTAS (APTAS) for a minimization problem is a family of
algorithms $\{A_\eps\}$ and a constant $c > 0$, such that there is an
algorithm $A_\eps$ for any $\eps > 0$, and $A_\eps$ outputs a
solution of value at most $(1 + \eps)\text{OPT} + c$, where OPT is
the value of an optimal solution. Suggest an APTAS for Bin Packing with Trees.

\subsection*{The algorithm}
\textbf{Solution outline} - First we will divide the items into large and small ones, the large ones into intervals of sizes, and we will round the sizes of the items in each interval to the maximum size in the interval. Then we will use the standard APTAS to get an APTAS for the abstracted sizes, we will resolve the conflicts, and at last we will add the small items to the bins.

\subsubsection*{Conflict resolution}
Now we will describe a conflict resolution algorithm for a tree conflict graph. The input $AbstractedBins$ is the solution produced by standard APTAS on the rounded instance: a list of bins, where each bin is a multiset of abstracted sizes (one slot per large item assigned to it). The goal is to assign each real item $i \in Large$ to a slot whose abstracted size equals $i$'s abstracted size, such that no two adjacent items share a bin.

The idea is to order the vertices so that at most one neighbor precedes each vertex. Given such an order, we greedily assign each vertex to an available slot of its abstracted size in a bin that contains no already-assigned neighbor. If no such bin exists, we open a new bin.

Here is an algorithm that finds the desired order of the vertices in the tree:
\begin{algorithm}
  \caption{1-inductive tree ordering}
  \begin{algorithmic}[1]
    \Procedure{OneInductiveTreeOrdering}{$T = (V, E)$}
      \State Pick any vertex $r \in V$ as the root
      \State \Return $\text{BFS}(T, r)$ \Comment{vertices in BFS discovery order}
    \EndProcedure
  \end{algorithmic}
\end{algorithm}

\begin{claim}
The order returned by the algorithm is a 1-inductive order.
\end{claim}

\begin{proof}
Since $T$ is a tree (acyclic), the BFS spanning tree is $T$ itself.
Every non-root vertex $v$ has exactly one parent, the vertex that
discovered it, which precedes $v$ in BFS order. Every other neighbor
of $v$ is a child of $v$ in the BFS tree, meaning $v$ discovered
it, so it appears \emph{after} $v$. The root has no preceding
neighbors at all. Hence every vertex has at most one preceding
neighbor.
\end{proof}

\begin{algorithm}
  \caption{Conflict resolution for Bin Packing with Trees}
  \begin{algorithmic}[1]
    \Procedure{ConflictResolution}{$AbstractedBins, T$}
    \Comment{$AbstractedBins$: list of bins with slots of abstracted sizes}
    \State $Order \gets$ \Call{OneInductiveTreeOrdering}{$T$}

    \State $ResolvedBins \gets []$
    \For{each vertex $v$ in $Order$}
      \State $\hat{s}_v \gets$ abstracted size of $v$
      \State Let $S$ be the set of bins in $AbstractedBins$ with an unmatched slot of size $\hat{s}_v$
      \State Filter $S$ to only include bins that do \emph{not} contain any already-matched neighbor of $v$

      \If{$S$ is not empty}
        \State Pick a bin $B \in S$, claim one slot of size $\hat{s}_v$ in $B$ for $v$, and add $B$ to $ResolvedBins$
      \Else
        \State Create a new bin $B'$, add $v$ to $B'$, and add $B'$ to $ResolvedBins$
      \EndIf
    \EndFor

    \State \Return $ResolvedBins$

    \EndProcedure
    \end{algorithmic}
\end{algorithm}

\textbf{Complexity analysis} - BFS runs in $O(n)$ time. We have at most $O(n)$ bins each containing at most $O(1/\eps)$ items. For each of the $n$ vertices, we can find and filter the available bins in $O(n)$ time. Thus, the overall complexity of the conflict resolution algorithm is $O(n^2)$.

\begin{claim}
The solution returned is conflict-free and feasible.
\end{claim}

\begin{proof}
By design, the algorithm never adds a vertex to a bin that already contains one of its neighbors, ensuring the solution is conflict-free. Furthermore, since rounding was upward, every item's real size satisfies $s_i \leq \hat{s}_i$. Therefore, the total real size of items placed in any bin from $AbstractedBins$ is at most the sum of their abstracted sizes, which is at most 1 (since $AbstractedBins$ is a feasible packing of the abstracted sizes). Any new bin created contains only a single item of size at most 1, so these are also feasible.
\end{proof}

\begin{claim}
The number of extra bins added is a constant depending only on $\eps$.
\end{claim}

\begin{proof}
Let $\hat{s}$ be a specific abstracted size. A new bin is opened for an item $v$ of size $\hat{s}$ only if all bins in $AbstractedBins$ with an unmatched slot of size $\hat{s}$ are filtered out. Since $v$ has at most one already-matched neighbor (from the 1-inductive order), it can filter out at most one bin. This implies that when a new bin is opened for size $\hat{s}$, all currently unmatched slots of size $\hat{s}$ must be contained in a single bin.

Because all large items have size at least $\eps$, any single bin can contain at most $\lfloor 1/\eps \rfloor$ slots. Therefore, at the moment the first new bin is opened for size $\hat{s}$, there are at most $\lfloor 1/\eps \rfloor$ unmatched slots of size $\hat{s}$ remaining in total.

Since the total number of items of size $\hat{s}$ equals the total number of slots of size $\hat{s}$ in $AbstractedBins$, every time we place an item of size $\hat{s}$ into a new bin, one slot of size $\hat{s}$ in $AbstractedBins$ will permanently remain empty. Consequently, the total number of new bins we can possibly open for size $\hat{s}$ is bounded by the number of available slots for size $\hat{s}$ when the first such new bin is created, which is at most $\lfloor 1/\eps \rfloor$.

There are at most $m = \lceil 1/\eps^2 \rceil$ distinct abstracted sizes. Thus, the total number of extra bins created across all sizes is at most $\lceil 1/\eps^2 \rceil \cdot \lfloor 1/\eps \rfloor \leq 1/\eps^3 + 1/\eps^2$, which is a constant depending only on $\eps$.
\end{proof}


\subsection*{Packing small items}
Since the conflict graph is a tree, it is bipartite, and we can 2-color the vertices. We will pack the small items greedily into the bins, and those that cannot be packed will be put in new bins according to their color.

\begin{algorithm}
  \caption{Packing small items for Bin Packing with Trees}
  \begin{algorithmic}[1]
    \Procedure{PackSmallItems}{$Bins, Small, T$}
      \State 2-color the vertices of $T$ with colors $1$ and $2$
      \State Let $Small_1, Small_2$ be the small items of colors $1$ and $2$
      \State Let $NewBins[1]$ and $NewBins[2]$ be empty bins

      \For{each $c \in \{1, 2\}$}
        \For{each $i \in Small_c$}
          \State $packed \gets \text{false}$
          \For{each $B \in Bins$}
            \State \textbf{if} $c = 2$ \textbf{and} $B \cap Small_1 \neq \emptyset$ \textbf{then continue}
            \If{sum of sizes in $B + s_i \leq 1$ \textbf{and} $B$ contains no neighbor of $i$}
              \State Add $i$ to $B$
              \State $packed \gets \text{true}$
              \State \textbf{break}
            \EndIf
          \EndFor
          \If{\textbf{not} $packed$}
            \If{sum of sizes in $NewBins[c] + s_i > 1$}
              \State Add $NewBins[c]$ to $Bins$
              \State $NewBins[c] \gets \text{empty bin}$
            \EndIf
            \State Add $i$ to $NewBins[c]$
          \EndIf
        \EndFor
      \EndFor

      \State Add any non-empty $NewBins[1]$ and $NewBins[2]$ to $Bins$

      \State \Return $Bins$
    \EndProcedure
  \end{algorithmic}
\end{algorithm}

\textbf{Complexity analysis} - 2-coloring the tree takes $O(n)$ time using BFS or DFS. The algorithm iterates over at most $n$ small items, and for each item, it checks at most $O(n)$ bins to see if it fits and has no conflicts. Checking for conflicts and capacity constraints can be done efficiently, leading to an overall running time of $O(n^2)$.

\begin{claim}
This algorithm yields a feasible packing of the small items.
\end{claim}

\begin{proof}
First, the capacity constraint is maintained: an item of size $s_i$ is only added to an existing bin if the sum of sizes including $s_i$ is at most 1. For the new color-specific bins, a fresh bin is opened if the capacity would exceed 1. Thus, no bin exceeds the size limit of 1.

Second, the conflict-free constraint is maintained. When adding to an existing bin, the algorithm explicitly verifies that the bin contains no neighbor of the item. When an item is placed in a new bin, it is grouped only with other items of the same color from the valid 2-coloring of the tree. Since adjacent vertices in a tree (which is bipartite) must have different colors, no two items of the same color can be neighbors. Therefore, the new bins are also strictly conflict-free. Thus, the resulting packing is feasible.
\end{proof}

\subsubsection*{Putting it all together}
Finally, we will describe the APTAS for Bin Packing with Trees.
\begin{algorithm}
  \caption{APTAS for Bin Packing with Trees}
  \begin{algorithmic}[1]
    \Procedure{BinPackingWithTrees}{$I, T, \eps$}
    % separate the items into small and large items
    \State $m = \lceil (\frac{2}{\eps})^2 \rceil$
    \State $Small \gets \{i \mid i \in I, s_i < \frac{\eps}{2}\}$
    \State $Large \gets I \setminus Small$
    % round the sizes of the large items
    \State sort $Large$ in non-increasing order of size
    \State Partition $Large$ into $m$ groups $G_1, G_2, \ldots, G_m$ of sizes as equal as possible
    \State For each $j \in [m]$ and $i \in G_j$, set $\hat{s}_i \gets \max_{i' \in G_j} s_{i'}$

    % get an APTAS for the abstracted sizes
    \State $Bins \gets$ standard APTAS($Large, \frac{\eps}{2}$)

    % resolve the conflicts
    \State $Bins \gets$ ConflictResolution($Bins, T$)

    % add the small items to the bins
    \State $Bins \gets$ PackSmallItems($Bins, Small, T$)

    \EndProcedure
  \end{algorithmic}
\end{algorithm}

\subsection*{Proof of correctness}
\begin{claim}
The solution returned is feasible.
\end{claim}

\begin{proof}
By the previous claims, each step of the algorithm maintains feasibility. The \textsc{ConflictResolution} algorithm ensures that all large items are packed into bins such that the sum of their actual sizes is at most 1, and no two adjacent items share a bin. The \textsc{PackSmallItems} algorithm guarantees that the addition of small items also respects both the capacity constraint and the conflict-free constraint. Since every item (both large and small) is placed into some bin without violating these constraints, the final solution is feasible.
\end{proof}

\subsection*{Running time analysis}
The algorithm's running time is dominated by the conflict resolution and small item packing steps, both of which take $O(n^2)$ time as analyzed earlier. The initial step of categorizing and rounding the items takes $O(n)$ time. The standard APTAS runs in polynomial time if $\eps$ is constant. Therefore, the overall running time of the algorithm is polynomial in $n$ for any fixed $\eps$.

\subsection*{Approximation ratio analysis}
\begin{claim}
For a given $\eps > 0$, the algorithm returns a packing into at most $(1 + \eps)\text{OPT} + c_\eps$ bins.
\end{claim}

\begin{proof}
Let $\text{OPT}$ be the optimal number of bins for the original instance.
The standard APTAS for the large items (with parameter $\eps/2$) produces a packing of the abstracted sizes into $m_{abstract} \leq (1 + \eps/2)\text{OPT} + c$ bins, where $c$ is a constant.

The \textsc{ConflictResolution} algorithm then produces a valid packing for the real large items, opening at most $c_{conflict} \leq 1/(\eps/2)^3 + 1/(\eps/2)^2$ additional bins, as proven earlier. Let $m_{large} = m_{abstract} + c_{conflict}$ be the total number of bins before packing small items. We have $m_{large} \leq (1 + \eps/2)\text{OPT} + c'$, where $c'$ is a constant depending only on $\eps$.

Next, consider the \textsc{PackSmallItems} phase. Let $U$ be the set of bins from the $m_{large}$ bins that are strictly less than $(1 - \eps/2)$ full. Any small item $i$ has size $s_i < \eps/2$, so it would fit into any bin in $U$ without violating the capacity constraint.

If an item $i$ is placed into a new color-specific bin, it must have been rejected by all bins in $U$. Since it fits by capacity, it must be rejected due to a conflict. This means $i$ has at least one neighbor in every bin in $U$.

Let $S'_c$ be the set of small items of color $c \in \{1, 2\}$ that are placed into new bins. Consider $S'_1$ first. When placing items of $Small_1$, no other small items have been packed yet. Thus, any conflict must be with a large item. Let $V_U$ be the set of large items packed in the bins of $U$. Since all large items have size at least $\eps/2$, each bin in $U$ contains at most $\lfloor 2/\eps \rfloor$ items. Thus, $|V_U| \leq |U| \cdot \lfloor 2/\eps \rfloor$.

The edges between $S'_1$ and $V_U$ form a bipartite subgraph of the tree $T$, which must be a forest. In a forest, the number of edges is strictly less than the number of vertices. Since every item in $S'_1$ has at least one large neighbor in every bin in $U$, there are at least $|S'_1| \cdot |U|$ edges. Therefore:
\[ |S'_1| \cdot |U| \leq |S'_1| + |V_U| - 1 \implies |S'_1|(|U| - 1) < |U| \cdot \lfloor 2/\eps \rfloor \]
The exact same logic applies to $S'_2$: because the algorithm explicitly forbids packing $Small_2$ into bins containing $Small_1$, any conflicts that force an item in $Small_2$ into a new bin must also be with large items in $V_U$. Thus, $|S'_2|(|U| - 1) < |U| \cdot \lfloor 2/\eps \rfloor$.

We consider two cases for $|U|$:
\begin{itemize}
    \item \textbf{Case 1: $|U| \geq 2$.} Then for each color $c$, $|S'_c| < \frac{|U|}{|U|-1} \lfloor 2/\eps \rfloor \leq 2 \lfloor 2/\eps \rfloor \leq 4/\eps$. This means at most $8/\eps$ small items are placed into new bins in total. Since each new bin contains items of the same color, and we only open a new bin when capacity is exceeded, these items will open at most $O(1/\eps)$ new bins. The total number of bins is $m_{large} + O(1/\eps) \leq (1 + \eps/2)\text{OPT} + c_\eps \leq (1 + \eps)\text{OPT} + c_\eps$.
    \item \textbf{Case 2: $|U| \leq 1$.} This means all but at most one of the $m_{large}$ bins are filled to at least $1 - \eps/2$. Let $m_{total}$ be the final number of bins. All newly opened bins (except possibly the last one for each color) are also filled to at least $1 - \eps/2$. Thus, at least $m_{total} - 3$ bins have a sum of sizes greater than $1 - \eps/2$. The total size of all items in the instance, which is a lower bound on OPT, is at least $(m_{total} - 3)(1 - \eps/2)$.
    Therefore, $m_{total} \leq \frac{\text{OPT}}{1 - \eps/2} + 3$. For $\eps \leq 1$, $\frac{1}{1 - \eps/2} \leq 1 + \eps$. So $m_{total} \leq (1 + \eps)\text{OPT} + 3$.
\end{itemize}
In both cases, the final number of bins is bounded by $(1 + \eps)\text{OPT} + c_\eps$ for some constant $c_\eps$ depending only on $\eps$, satisfying the definition of an APTAS.
\end{proof}
